% chapter1.tex -- en (English)
\chapter{Preface}

Thank you for interesting in {\autor}'s {\Bezeichnung}.\\
This is a control unit for PCs, which can emulate keyboard and mouse
actions and is inspired on the gaming controllers of various game 
consoles. It is based on the Arduino board \textit{Esplora} which is
technically a \textit{Arduino Leonardo} and it already contains some
sensors and actuators. The printed circuit board was designed roughly in
the form of a gaming controller.

\url{https://store.arduino.cc/arduino-esplora}

\section{Legal Notes}
When designing the {\Bezeichnung}, it was taken care to use only 
software that is provided under a free license such as GNU GPL or 
similar or that has been released into the public domain.

\subsection*{Trademarks}
Some names used in this document may be trademarks. The use of these 
trademarks by third parties for their purposes may infringe the rights 
of the holders.

\subsection*{Links}
This manual contains links to external sites on the internet. Despite
linking the author {\autor} does not appropriate these contents, since
they are not within his sphere of influence! At the time of linking
there were no illegal contents noticeable. It is not reasonable for the
author to check the links permanently for any changes that might violate
the law. However, if current or future content should be illegal, the 
author may be contacted by e-mail to
\url{mailto:himself@schlizbaeda.de}.
Appropriate actions will then be taken to remove the affected link(s).

\subsection*{Licensing of this project's components}
\uline{Code:}\\
Most Arduino code examples are in the public domain and may be used
by any person. The code for {\Bezeichnung} was developed from the
templates\\
\url{https://www.arduino.cc/en/Tutorial/EsploraJoystickMouse},\\
\url{https://www.arduino.cc/en/Tutorial/EsploraKart} and\\
\url{https://github.com/circuit69/EsploraTinkerkit}\\
Due to many adjustments and improvements the author decided to publish
his work under the GNU GPL v3:\\
\url{https://www.gnu.org/licenses/gpl-3.0.en.html}\\
\includegraphics[height=35px]{GPLv3.png}

\uline{Housing:}\\
A CAD file collection for 3D printing was found at
\url{https://www.thingiverse.com/thing:45880}. This data is published
under CC BY-NC 3.0
(\url{http://creativecommons.org/licenses/by-nc/3.0/}). 
This license does not allow commercial use.\\
\includegraphics[height=35px]{CC-BY-NC.png}


\subsection*{Image Copyright}
All pictures, photographs and technical drawings were created by the
author himself. They are published under the \textit{Creative Commons}
license \textbf{CC BY-SA 3.0}. Therefore anyone may use these images
in original form or adapted by giving name attribution of the author.\\
\includegraphics[height=35px]{CC-BY-SA.png}

The Arduino logo is in the public domain:\\
\includegraphics[height=35px]{arduino.png}


\newpage
\section{Remarks to a genderwise correctly written text}
Grammatical gender differentiation is very unusual in English language
compared to German language. In  German speaking countries there is
currently a trend of \textit{political correctness} to mix the
grammatical gender with the biological sex. This results in
\textit{genderwise correctly} written texts which are sometimes really
difficult to read.\\
This translation isn't affected by this German gender trend anyway. So
this may be regarded as a typical German problem \smiley{wink}

\section{Achnowledgements}
{\autor} would like to welcome the user \textbf{@dale}
(\url{https://forum-raspberrypi.de/user/23726-dale/}) from the German
Raspberry Pi forum (\url{https://forum-raspberrypi.de}):\\
He printed the housing on his 3D printer.

\section{Abstract} % Brief Description
The {\Bezeichnung} is a control unit with USB2.0 connection inspired by
the game console based gaming controllers. Instead of proprietary
signals this device transmits keyboard and mouse commands when the user
actuates the switches of the device.

Due to its open-source and free programming, it can be adapted to many
applications in the source code by assigning the desired keyboard codes
or mouse commands to the control elements (sensors) of the device.

There are eight different assignments (so-called modes) for several
applications which can be switched during operation. The selected 
assignment is indicated by the colour of the three-colour LED on the
Esplora board.

The software of the {\Bezeichnung} emulates a USB keyboard and a USB
mouse and can be used immediately on most of the PC operating systems
like \textit{Windows} or \textit{GNU/Linux} as well as on a {\RPi} under
\textit{Raspbian} without any additional driver installation.
